\section{Single-round setup}
\label{sec:setup}

\paragraph{Model, task, and training.}
Following \citet{shin2026compositional}---a compositional variant of
self-training on a model's own outputs \citep{lee2025selfimproving}---we train
a character-level Transformer (approximately 14M parameters, no positional
encoding) from scratch to add two numbers of equal length. A starting model is
first trained with correct answers on sums of 3 to 7 digits, which we call
\emph{easy problems}. Four successive rounds then extend the maximum length to
9, 11, 13, and 15 digits. In each round, the teacher generates answers to fresh
problems at every length from 8 digits up to the current maximum until it has
kept 10,000 answers per length, and the model trains on these together with
easy problems that it answers itself. Sums are split into two pieces up to 12
digits and three at 13 to 15, so every piece has 3 to 7 digits. Each length
uses three fixed splits (five at 13 to 15 digits) with no shared cut positions;
the first supplies the candidate answer, and the others serve the agreement
check and the agreement rule.

\paragraph{Teachers.}
With the no-carry filter of \citet{shin2026compositional}, every addition round
in our reproduction paid off. Testing a three-way predictor requires rounds
that also stall or degrade the model, so we construct a benchmark spanning all
three outcomes by design. We run four teachers: the \emph{no-carry rule} and
\emph{agreement rule} (Section~\ref{sec:two-roles}), expected to help; \emph{no
selection}, which keeps the first split's answer unconditionally and thus
trains on every carry error; and \emph{sample and vote}, which does not split
at all, instead sampling five direct answers and keeping one that appears at
least three times \citep{wang2023selfconsistency}. Additional cases come from
starting models stopped early (weak piece accuracy), from applying \emph{no
selection} to a model previously trained with the no-carry rule (so accuracy
can fall from near 100\%), and from retraining lengths already learned in
earlier rounds.

\paragraph{Cases and outcomes.}
A \emph{case} is one combination of teacher, starting model, round, and problem
length. Before the round, the checks run on 2,000 fresh problems of that length
to issue a prediction; afterwards, the outcome is scored from held-out test
accuracy before and after training. At a length the model already knew, a
change within 3 points counts as \emph{did nothing}, and a larger drop---at
that length or on easy problems---counts as \emph{did harm}. At a new length,
accuracy starts near zero and cannot fall, nor does a rise alone show that the
right rule was learned, so we compare against the best accuracy achieved by any
teacher at that length in the same round and from the same starting model: a
rise of at least 20 points ending within 10 points of the best \emph{paid off};
a rise ending 30 points or more below the best taught wrong answers and counts
as \emph{did harm}; a rise under 10 points \emph{did nothing}. A few
intermediate cases (such as partially learned lengths that rise further) are
left unclassified by this rule and examined separately in
Section~\ref{sec:errors}.

\paragraph{Calibration and evaluation.}
We trained five starting models from distinct random seeds: two for the
calibration runs used to set the rule of Section~\ref{sec:rule}, and three held
out for evaluation. Before any evaluation run began, the prediction rule, the
outcome rule, and a pre-specified pass mark---at least 90\% accuracy overall,
with at least five cases each of \emph{did nothing} and \emph{did harm}---were
frozen in a tagged commit. The three evaluation models produce 270 cases, of
which the outcome rule classifies 258. Because not every analysis in this paper
was fixed at the same moment, Appendix~\ref{app:fixed-when} records when
each was specified (Table~\ref{tab:fixed-when}).
